\documentclass[11pt,a4paper]{article}

\usepackage[T1]{fontenc}
\usepackage[utf8]{inputenc}
\usepackage[italian]{babel}
\usepackage{lmodern}
\usepackage{microtype}
\usepackage[a4paper,margin=2.15cm,headheight=15pt]{geometry}
\usepackage{booktabs,longtable,tabularx,array}
\usepackage{ragged2e}
\usepackage{enumitem}
\usepackage{xcolor}
\usepackage[most]{tcolorbox}
\usepackage{fancyhdr}
\usepackage{hyperref}
\usepackage{xurl}
\usepackage{listings}
\usepackage{titlesec}
\usepackage{amsmath,amssymb}

\definecolor{ddtadark}{RGB}{28,38,52}
\definecolor{ddtagray}{RGB}{244,246,248}
\definecolor{ddtablue}{RGB}{42,88,135}
\definecolor{ddtagreen}{RGB}{48,111,78}
\definecolor{ddtaorange}{RGB}{148,91,25}
\definecolor{ddtared}{RGB}{140,48,48}
\definecolor{ddtapurple}{RGB}{92,63,122}

\hypersetup{
  colorlinks=true,
  linkcolor=ddtablue,
  urlcolor=ddtablue,
  citecolor=ddtablue,
  pdftitle={Base Analysis Operational Guide},
  pdfauthor={Documentation-Driven Threat Analysis research workspace}
}

\pagestyle{fancy}
\fancyhf{}
\lhead{DDTA -- Base Analysis Operational Guide}
\rhead{Revision 3 - R25 pre-holdout}
\cfoot{\thepage}

\titleformat{\section}{\Large\bfseries\color{ddtadark}}{\thesection}{0.7em}{}
\titleformat{\subsection}{\large\bfseries\color{ddtadark}}{\thesubsection}{0.7em}{}
\titleformat{\subsubsection}{\normalsize\bfseries\color{ddtadark}}{\thesubsubsection}{0.7em}{}

\setlist[itemize]{leftmargin=1.55em,itemsep=0.22em,topsep=0.35em}
\setlist[enumerate]{leftmargin=1.75em,itemsep=0.22em,topsep=0.35em}
\setlength{\parskip}{0.35em}
\setlength{\parindent}{0pt}
\setlength{\emergencystretch}{3em}

\lstset{
  basicstyle=\ttfamily\small,
  columns=fullflexible,
  keepspaces=true,
  frame=single,
  framerule=0.3pt,
  rulecolor=\color{ddtablue},
  backgroundcolor=\color{ddtagray},
  xleftmargin=0.6em,
  xrightmargin=0.6em,
  aboveskip=0.5em,
  belowskip=0.5em,
  breaklines=true,
  showstringspaces=false
}

\newtcolorbox{closedbox}[1][REGOLA STABILE]{enhanced,breakable,colback=ddtagray,colframe=ddtagreen,title={#1},fonttitle=\bfseries}
\newtcolorbox{workingbox}[1][REGOLA BASELINE-SCOPED]{enhanced,breakable,colback=white,colframe=ddtaorange,title={#1},fonttitle=\bfseries}
\newtcolorbox{openbox}[1][NOT SPECIFIED / DIAGNOSTIC]{enhanced,breakable,colback=white,colframe=ddtapurple,title={#1},fonttitle=\bfseries}
\newtcolorbox{goodbox}[1][PREFERIBILE]{enhanced,breakable,colback=white,colframe=ddtagreen,title={#1},fonttitle=\bfseries}
\newtcolorbox{badbox}[1][DA EVITARE]{enhanced,breakable,colback=white,colframe=ddtared,title={#1},fonttitle=\bfseries}
\newtcolorbox{examplebox}[1][ESEMPIO FACIAL ACCESS]{enhanced,breakable,colback=ddtagray,colframe=ddtablue,title={#1},fonttitle=\bfseries}
\newtcolorbox{conceptbox}[1][CONCETTO]{enhanced,breakable,colback=white,colframe=ddtablue,title={#1},fonttitle=\bfseries}

\newcolumntype{Y}{>{\RaggedRight\arraybackslash}X}
\newcolumntype{L}[1]{>{\RaggedRight\arraybackslash}p{#1}}

\newcommand{\code}[1]{\nolinkurl{#1}}
\newcommand{\BA}{\textit{Base Analysis}}
\newcommand{\BAR}{\textit{BAReferent}}
\newcommand{\BAP}{\textit{BAProposition}}

\begin{document}

\begin{titlepage}
\centering
\vspace*{1.5cm}
{\Large Documentation-Driven Threat Analysis\par}
\vspace{0.8cm}
{\Huge\bfseries Base Analysis Operational Guide\par}
\vspace{0.25cm}
{\LARGE Revision 3 - R25 pre-holdout BA contract\par}
\vspace{0.9cm}
{\large Procedura source-first dalla baseline governata alla accepted Base Analysis e alle projection downstream\par}
\vfill
\begin{minipage}{0.88\textwidth}
\small
\textbf{Stato.} R25 PHASE 2 COMPLETE / FROZEN PRE-HOLDOUT BA OPERATIONAL GUIDE.

\medskip
\textbf{Continuation baseline.} \code{7209f477d7de3b8d8726c91ce93682dddfddb38a}.

\medskip
\textbf{Pre-holdout BA contract set.} BA0 R1; BA1 R1; BA2 R3; BA3 R1; BA4 R1; BA5 R1.

\medskip
\textbf{Project authority example.} \code{FACIAL-ACCESS-GOV-R2}.

\medskip
\textbf{Facial Access BA evidence.} \code{FACIAL-ACCESS-BA-R24-R1}: complete case BA / post-BA regression PASS.

\medskip
\textbf{Integrated BA6 acceptance.} OPEN.

\medskip
Questa guida non sostituisce i contratti normativi BA0--BA5 né i documenti governati del progetto. Li organizza in una procedura leggibile e operativa. BA6 resta il successivo gate di validazione integrata, non un risultato già chiuso.
\end{minipage}
\vfill
{\large 29 agosto 2026\par}
\end{titlepage}

\pagenumbering{roman}
\tableofcontents
\newpage
\pagenumbering{arabic}

\section{Scopo e stato pre-holdout}

Questa revisione consolida la procedura operativa della Base Analysis prima della selezione del nuovo holdout.

\begin{closedbox}[Pre-holdout boundary]
Il contratto che il futuro holdout potrà falsificare è:
\begin{lstlisting}
BA0 R1
BA1 R1
BA2 R3
BA3 R1
BA4 R1
BA5 R1
\end{lstlisting}
BA6 non appartiene a questo risultato congelabile: resta il gate integrato da valutare dopo il nuovo caso, le projection multiple e il change/revalidation cycle.
\end{closedbox}

Al momento di questa revisione:

\begin{lstlisting}
holdout domain
    NOT SELECTED

holdout documentation
    NOT AUTHORED

BA6 integrated acceptance
    OPEN
\end{lstlisting}

\subsection{Anti-contamination rule}

Il futuro holdout non deve essere scelto o scritto per esercitare operatori, ruoli o strutture BA specifiche. La documentazione del caso deve essere governata dal problema e dalle decisioni del progetto; la BA deve poi tentare di rappresentarla source-first.

Un failure futuro può riaprire soltanto il minimo owning layer attraverso un counterexample concreto.

\section{Che cos'è la Base Analysis}

Per una baseline documentale governata, la Base Analysis è la rappresentazione analitica condivisa, methodology-neutral, del significato di progetto necessario a DDTA.

\begin{closedbox}[Authority boundary]
\begin{lstlisting}
governed project documentation
    = project authority

accepted Base Analysis
    = shared analytical baseline

projection / threat method / tests
    = downstream consumers
\end{lstlisting}
\end{closedbox}

La BA può normalizzare e rendere riusabile il significato, ma non può creare silenziosamente project commitment. Un finding downstream può localizzare una review necessaria; la project truth cambia soltanto attraverso documentazione governata.

\section{Il contratto BA0--BA5 e il gate BA6}

La metodologia separa responsabilità distinte.

\begin{longtable}{L{0.13\textwidth}L{0.75\textwidth}}
\toprule
\textbf{Layer} & \textbf{Responsabilità} \\
\midrule\endhead
BA0 & authority boundary, minimalità, unresolved state, feedback direction. \\
BA1 & due sole famiglie first-class: \code{BAReferent} e \code{BAProposition}. \\
BA2 R3 & proposition semantics, operatori, ruoli, cardinalità e strutture locali. \\
BA3 & provenance, derivazione, review/freshness, continuity e revalidation. \\
BA4 & projection, traceability, interpretation e coverage downstream. \\
BA5 & canonical semantic registries e controlled authoring. \\
BA6 & successivo gate di accettazione integrata; attualmente OPEN. \\
\bottomrule
\end{longtable}

BA6 non rende BA0--BA5 ``più veri'': verifica successivamente che il metodo funzioni come catena integrata sotto un validation protocol più forte.

\section{Regola di allineamento del contract set}

BA1, BA3, BA4 e BA5 sono artifact storicamente chiusi. Alcune loro note di allineamento o snapshot descrivono lo stato BA2 esistente al momento della rispettiva chiusura.

Per l'esecuzione R25 corrente vale:

\begin{lstlisting}
current BA2 operator semantics
    -> BA2_RELATION_ACTION_VOCABULARY_R3

current BA2 operator-role semantics
    -> BA2_RELATION_ACTION_VOCABULARY_R3

BA1 / BA3 / BA4 / BA5 normative core
    -> retained

older forward pointer / embedded registry snapshot
    -> historical alignment context
    -> not current BA2 authority
\end{lstlisting}

In particolare, lo snapshot BA5 a tredici operatori precede \code{decisionRule}. BA5 continua a governare canonicalità, domain-scoped resolution, revision identity ed extension workflow; l'elenco operativo corrente degli operatori e dei ruoli deriva da BA2 R3.

Questa regola evita di riscrivere closed contract bodies soltanto per aggiornare cronologia o pointer.

\section{BAReferent e BAProposition}

\subsection{BAReferent}

Un \BAR{} è un significato di progetto methodology-neutral con identità indipendente quando deve essere:

\begin{itemize}
\item riusato da più proposition;
\item qualificato o vincolato;
\item correlato con altri significati;
\item confrontato tra baseline;
\item target di più asserzioni;
\item tracciato stabilmente attraverso change/revalidation.
\end{itemize}

Può denotare participant, capability, information, behavior, state, service, contract, boundary o altro significato. Queste categorie non sono automaticamente nuove famiglie BA.

\subsection{BAProposition}

Una \BAP{} è una asserzione analitica su uno o più BAReferent.

\begin{closedbox}[No proposition proxy]
La proposition non è il comportamento o l'oggetto di progetto che descrive. Se quel significato deve essere riusato indipendentemente, riceve BAReferent identity.
\end{closedbox}

Questa regola giustifica, nel caso Facial Access, \code{RecognitionCaptureDelivery} come BAReferent separato dalla proposition \code{transfer} che lo descrive.

\section{Il criterio di dettaglio}

La BA non replica l'intera documentazione e non ricostruisce automaticamente una architettura completa.

\begin{lstlisting}
more detail
    only if needed to:
      preserve a governed distinction
      answer a material downstream question
      or expose a missing governed fact
\end{lstlisting}

Stopping test:

\begin{conceptbox}
Abbiamo abbastanza semantica per formulare la domanda e distinguere evidenza supportata, contraddetta o \code{NOT SPECIFIED}?
\end{conceptbox}

Se sì, stop. Se no, si identifica la distinzione mancante o si produce una diagnosis. Non si inventa l'implementazione.

\section{Procedura operativa source-first}

\begin{enumerate}
\item Pin della project-document authority e della source revision.
\item Lettura source-first del governed source set.
\item Identificazione del minimum shared meaning necessario.
\item Creazione dei BAReferent richiesti dall'identità riusabile.
\item Formulazione delle BAProposition usando BA2 R3.
\item Collegamento BA3 alla source evidence e agli eventuali derivation/revalidation contexts.
\item Registrazione esplicita di ambiguity, diagnostic e \code{NOT SPECIFIED}.
\item Acceptance preliminare del minimum BA per lo scope dichiarato.
\item Regression completa verso le fonti e forbidden-inference review.
\item Pressure review della metodologia soltanto su counterexample concreti.
\item Materializzazione/revisione della accepted case BA.
\item Projection/downstream consumption secondo BA4.
\item Se il validation protocol lo richiede, governed change, BA3 revalidation, rebuild e re-analysis.
\item BA6 integrated verdict soltanto dopo l'intero gate previsto.
\end{enumerate}


\section{Dalla documentazione alla BA: come si estrae il significato}

Questa è la regola pratica centrale della guida. Non si parte da un operatore BA cercando una frase che lo giustifichi. Si parte dalla documentazione governata e si procede sempre nello stesso ordine:

\begin{lstlisting}
1. documentazione di partenza
2. significati che la documentazione governa davvero
3. BAReferent che devono avere identita' riusabile
4. BAProposition minima che preserva quei significati
5. significati che NON sono governati e non devono essere inventati
\end{lstlisting}

\begin{closedbox}[Regola di estrazione]
Il verbo naturale presente nel requisito è un indizio, non il nome automatico dell'operatore BA. L'operatore viene scelto soltanto dopo aver stabilito quale distinzione semantica la documentazione governa.
\end{closedbox}

Gli esempi marcati \textbf{Facial Access} derivano dalla documentazione governata corrente e mostrano il tipo di lettura usato nella Base Analysis del caso. Gli esempi marcati \textbf{DIDATTICO NON NORMATIVO} sono piccoli frammenti inventati soltanto per spiegare operatori che il caso Facial Access non esercita in modo sufficientemente chiaro.

\subsection{Esempio base: da FR-3.4.2 a BAReferent e BAProposition}

\begin{examplebox}[DOCUMENTAZIONE DI PARTENZA -- Facial Access FR-3.4.2]
Per ogni \code{RecognitionCapture} destinata al riconoscimento da parte di \code{RecognitionProcessor}, \code{CameraSubsystem} deve consegnarla a \code{RecognitionProcessor} mantenendo il binding con la stessa \code{IdentityDeterminationRequest} per cui la capture è stata acquisita.
\end{examplebox}

\textbf{Passo 1 -- significati governati.}

La frase governa almeno cinque significati distinti:

\begin{itemize}
\item esiste un soggetto che effettua la consegna: \code{CameraSubsystem};
\item esiste un destinatario: \code{RecognitionProcessor};
\item esiste un contenuto trasferito: \code{RecognitionCapture};
\item esiste il comportamento di delivery, perché altre Decision e Security Requirement lo qualificano;
\item la capture resta legata alla stessa \code{IdentityDeterminationRequest}.
\end{itemize}

\textbf{Passo 2 -- BAReferent.}

Servono quindi identità riusabili per:

\begin{lstlisting}
CameraSubsystem
RecognitionProcessor
RecognitionCapture
RecognitionCaptureDelivery
IdentityDeterminationRequest
\end{lstlisting}

\code{RecognitionCaptureDelivery} non viene creato perché ``ogni verbo deve diventare un oggetto''. Riceve una propria identità perché il comportamento di delivery è successivamente target del consumed service, del medium e delle security property.

\textbf{Passo 3 -- BAProposition.}

La consegna diventa:

\begin{lstlisting}
transfer
  behavior    -> RecognitionCaptureDelivery
  source      -> CameraSubsystem
  destination -> RecognitionProcessor
  content     -> RecognitionCapture
\end{lstlisting}

Il binding alla richiesta è una distinzione diversa e viene preservato separatamente:

\begin{lstlisting}
correlate
  correlatedItem     -> RecognitionCapture
  correlationContext -> IdentityDeterminationRequest
\end{lstlisting}

\textbf{Passo 4 -- cosa non estrarre.}

Da quella frase non possiamo aggiungere:

\begin{lstlisting}
protocol = HTTPS
acknowledgement = REQUIRED
retry = 3
transport ownership = CameraSubsystem
RecognitionCaptureDelivery = successful
\end{lstlisting}

perché nessuno di questi fatti è governato da FR-3.4.2.

\subsection{Come nasce un BAReferent}

Un nome nella documentazione non diventa automaticamente BAReferent. La domanda è: \emph{questo significato deve essere riusato, qualificato, correlato o tracciato come identità indipendente?}

\begin{examplebox}[DOCUMENTAZIONE DI PARTENZA -- Facial Access]
La documentazione distingue \code{AccessAuthorizationState} dalla relativa \code{GovernedIdentity}; lo stato riguarda l'identità ma non costituisce né crea l'identità stessa.
\end{examplebox}

Qui entrambi i significati devono essere richiamabili indipendentemente. Sono quindi BAReferent distinti:

\begin{lstlisting}
AccessAuthorizationState [BAReferent]
GovernedIdentity         [BAReferent]
\end{lstlisting}

La relazione tra loro verrà espressa da una proposition. Non si fonde invece tutto in un unico oggetto come:

\begin{lstlisting}
GovernedIdentityWithAuthorizationState
\end{lstlisting}

perché la documentazione governa proprio la distinzione tra identità e stato.

\subsection{Come nasce una BAProposition}

Una BAProposition rappresenta una asserzione governata tra significati. Non è un duplicato della frase sorgente e non è un contenitore arbitrario di prose.

\begin{examplebox}[DOCUMENTAZIONE DI PARTENZA -- Facial Access D-3.5]
La realizzazione di MR-0003 consuma un servizio di connettività disponibile per supportare il delivery della \code{RecognitionCapture} e non possiede né gestisce l'infrastruttura di trasporto sottostante.
\end{examplebox}

I significati riusabili sono almeno:

\begin{lstlisting}
IdentityDeterminationRealization
ConnectivityService
\end{lstlisting}

L'asserzione minima è:

\begin{lstlisting}
consumeService
  consumer -> IdentityDeterminationRealization
  service  -> ConnectivityService
\end{lstlisting}

La negazione dell'ownership rimane una boundary da preservare; non viene invertita creando una relazione di ownership verso il consumer.

\section{I quattordici operatori letti dalla documentazione}

Questa sezione mostra, per ogni operatore BA2 R3, quale tipo di documentazione lo giustifica e come si arriva alla proposition.

\subsection{transfer -- qualcosa viene trasferito da una source a una destination}

\begin{examplebox}[DOCUMENTAZIONE DI PARTENZA -- Facial Access FR-3.4.2]
\code{CameraSubsystem} deve consegnare la \code{RecognitionCapture} a \code{RecognitionProcessor}.
\end{examplebox}

\textbf{Lettura semantica:} la documentazione governa una conveyance con origine, destinazione e contenuto.

\begin{lstlisting}
BAReferent:
  CameraSubsystem
  RecognitionProcessor
  RecognitionCapture
  RecognitionCaptureDelivery   # necessario perche' il delivery e' qualificato altrove

BAProposition:
  transfer
    behavior    -> RecognitionCaptureDelivery
    source      -> CameraSubsystem
    destination -> RecognitionProcessor
    content     -> RecognitionCapture
\end{lstlisting}

\textbf{Non estrarre:} protocollo, hop, acknowledgement, retry, cifratura o successo della consegna se la documentazione non li governa.

\subsection{produce -- un actor rende disponibile un risultato governato}

\begin{examplebox}[DOCUMENTAZIONE DI PARTENZA -- Facial Access FR-3.2.1]
La capability di riconoscimento deve produrre un esito governato della determinazione dell'identità distinguendo determinazione riuscita, negativa e non conclusiva.
\end{examplebox}

D-3.4 assegna l'esecuzione del riconoscimento a \code{RecognitionProcessor}. La lettura combinata permette di mantenere actor e result separati:

\begin{lstlisting}
BAReferent:
  RecognitionProcessor
  IdentityDeterminationOutcome

BAProposition:
  produce
    actor  -> RecognitionProcessor
    input  -> RecognitionCapture
    result -> IdentityDeterminationOutcome
\end{lstlisting}

\textbf{Non estrarre:} score, confidence, threshold, ranking o algoritmo di decisione, esplicitamente non governati.

\subsection{create -- viene stabilito un nuovo item o evento}

\begin{examplebox}[DOCUMENTAZIONE DI PARTENZA -- DIDATTICO NON NORMATIVO]
Quando una \code{RegistrationRequest} viene approvata, \code{AccountProvisioning} deve creare un nuovo \code{UserAccount} e associarlo a quella richiesta.
\end{examplebox}

La parola importante non è soltanto ``creare'': la frase governa che dopo l'azione esiste un nuovo account che prima non esisteva come item di progetto.

\begin{lstlisting}
BAReferent:
  AccountProvisioning
  RegistrationRequest
  UserAccount

BAProposition:
  create
    actor   -> AccountProvisioning
    created -> UserAccount

  correlate
    correlatedItem     -> UserAccount
    correlationContext -> RegistrationRequest
\end{lstlisting}

\textbf{Non estrarre:} schema database, ID generation, password iniziale o persistence technology.

\subsection{observe -- un actor legge o ispeziona qualcosa che esiste}

\begin{examplebox}[DOCUMENTAZIONE DI PARTENZA -- DIDATTICO NON NORMATIVO]
\code{MonitoringAgent} deve leggere il \code{CurrentDoorState} esposto da \code{DoorController} senza modificarlo.
\end{examplebox}

Qui il requisito governa una lettura di significato esistente e nega implicitamente che quella operazione sia la creazione o la transizione dello stato.

\begin{lstlisting}
BAReferent:
  MonitoringAgent
  CurrentDoorState

BAProposition:
  observe
    actor    -> MonitoringAgent
    observed -> CurrentDoorState
\end{lstlisting}

\textbf{Non estrarre:} polling interval, API, cache o un nuovo \code{ObservationRecord} se non sono documentati.

\subsection{transition -- cambia lo stato di un subject}

\begin{examplebox}[DOCUMENTAZIONE DI PARTENZA -- DIDATTICO NON NORMATIVO]
Quando \code{AccountAdministration} revoca una credenziale attiva, \code{AccessCredential} deve passare dallo stato \code{ACTIVE} allo stato \code{REVOKED}.
\end{examplebox}

\begin{lstlisting}
BAReferent:
  AccountAdministration
  AccessCredential
  ActiveCredentialState
  RevokedCredentialState

BAProposition:
  transition
    actor     -> AccountAdministration
    subject   -> AccessCredential
    fromState -> ActiveCredentialState
    toState   -> RevokedCredentialState
\end{lstlisting}

\textbf{Non estrarre:} timestamp, persistence, rollback o workflow intermedio se non governati.

\subsection{correlate -- più significati devono appartenere allo stesso contesto}

\begin{examplebox}[DOCUMENTAZIONE DI PARTENZA -- Facial Access FR-3.4.2]
La \code{RecognitionCapture} deve mantenere il binding con la stessa \code{IdentityDeterminationRequest} per cui è stata acquisita.
\end{examplebox}

La parola decisiva è ``stessa'': non basta sapere che la capture fa riferimento a qualche request. La documentazione governa identità di contesto.

\begin{lstlisting}
BAReferent:
  RecognitionCapture
  IdentityDeterminationRequest

BAProposition:
  correlate
    correlatedItem     -> RecognitionCapture
    correlationContext -> IdentityDeterminationRequest
\end{lstlisting}

\textbf{Non sostituire con} \code{reference}: perderebbe la garanzia di same-request binding.

\subsection{reference -- A riguarda o richiama B senza semantica più forte}

\begin{examplebox}[DOCUMENTAZIONE DI PARTENZA -- Facial Access MR-0002]
\code{AccessAuthorizationState} riguarda una specifica \code{GovernedIdentity}; i due significati restano distinti e lo stato non crea l'identità.
\end{examplebox}

\begin{lstlisting}
BAReferent:
  AccessAuthorizationState
  GovernedIdentity

BAProposition:
  reference
    referencer -> AccessAuthorizationState
    referenced -> GovernedIdentity
\end{lstlisting}

\textbf{Non estrarre:}

\begin{lstlisting}
AccessAuthorizationState.authorized = TRUE | FALSE
ownership(AccessAuthorizationState, GovernedIdentity)
\end{lstlisting}

perché la documentazione non governa quel vocabolario né quelle relazioni.

\subsection{dependOn -- un significato richiede un prerequisite}

\begin{examplebox}[DOCUMENTAZIONE DI PARTENZA -- Facial Access D-3.4]
La \code{RecognitionCapture} deve essere resa disponibile da \code{CameraSubsystem} a \code{RecognitionProcessor}; l'indisponibilità di tale interazione può impedire il completamento della determinazione dell'identità.
\end{examplebox}

Questa conseguenza governa una dipendenza funzionale: il completamento della realizzazione di identity determination richiede che il delivery necessario sia disponibile.

\begin{lstlisting}
BAReferent:
  IdentityDeterminationOutcome
  RecognitionCaptureDelivery

BAProposition:
  dependOn
    dependent    -> IdentityDeterminationOutcome
    prerequisite -> RecognitionCaptureDelivery
\end{lstlisting}

\textbf{Non estrarre:} ownership, provider, protocollo o causa tecnica dell'indisponibilità.

\subsection{consumeService -- il progetto usa un servizio senza acquisirne ownership}

\begin{examplebox}[DOCUMENTAZIONE DI PARTENZA -- Facial Access D-3.5]
La realizzazione di MR-0003 consuma un servizio di connettività disponibile e non possiede né gestisce l'infrastruttura di trasporto sottostante.
\end{examplebox}

\begin{lstlisting}
BAReferent:
  IdentityDeterminationRealization
  ConnectivityService

BAProposition:
  consumeService
    consumer -> IdentityDeterminationRealization
    service  -> ConnectivityService
\end{lstlisting}

Il ruolo opzionale \code{provider} resta assente finché la documentazione non identifica un provider governato.

\textbf{Non estrarre:} provider implicito, ownership, gestione dell'infrastruttura o proprietà del servizio non garantite.

\subsection{realize -- un elemento concreto realizza un significato più astratto}

\begin{examplebox}[DOCUMENTAZIONE DI PARTENZA -- Facial Access D-3.1]
Il progetto usa riconoscimento facciale come strategia corrente per determinare quale identità governata corrisponde alla persona presente al punto di accesso.
\end{examplebox}

La documentazione distingue la responsabilità astratta di identity determination dalla strategia concreta scelta nella baseline.

\begin{lstlisting}
BAReferent:
  IdentityDetermination
  FacialRecognitionStrategy

BAProposition:
  realize
    abstract    -> IdentityDetermination
    realization -> FacialRecognitionStrategy
\end{lstlisting}

\textbf{Non estrarre:} modello ML, sensore, score, threshold, subclassing o deployment topology. \code{realize} dice soltanto quale significato concreto materializza quello astratto nel contesto governato.

\subsection{assignResponsibility -- la documentazione colloca una responsabilità}

\begin{examplebox}[DOCUMENTAZIONE DI PARTENZA -- DIDATTICO NON NORMATIVO]
\code{SecurityOperationsTeam} è responsabile dell'approvazione delle richieste di accesso di emergenza.
\end{examplebox}

Per rendere la responsibility interrogabile senza inventare un ruolo organizzativo generico:

\begin{lstlisting}
BAReferent:
  SecurityOperationsTeam
  EmergencyAccessApproval
  ApprovalResponsibility

BAProposition:
  assignResponsibility
    responsibleParty    -> SecurityOperationsTeam
    responsibilityScope -> EmergencyAccessApproval
    responsibilityKind  -> ApprovalResponsibility
\end{lstlisting}

\textbf{Non estrarre:} ownership di sistemi, autorizzazioni tecniche o separazione dei compiti se non documentate.

\subsection{constrain -- un target è limitato da un valore o dominio governato}

\begin{examplebox}[DOCUMENTAZIONE DI PARTENZA -- Facial Access D-3.6]
\code{CameraSubsystem} e \code{RecognitionProcessor} utilizzano connettività Ethernet cablata per il delivery della \code{RecognitionCapture}.
\end{examplebox}

La proprietà è locale al comportamento di delivery, non all'intero sistema né al servizio consumato end-to-end.

\begin{lstlisting}
BAReferent:
  RecognitionCaptureDelivery

BAProposition:
  constrain
    constraintTarget -> RecognitionCaptureDelivery
    constraintValue
      property   -> interconnectionMedium
      vocabulary -> [WIRED_ETHERNET]
\end{lstlisting}

\textbf{Non estrarre:}

\begin{lstlisting}
ConnectivityService.medium = WIRED_ETHERNET end-to-end
protocol = Ethernet/IP/TCP
network topology = ...
\end{lstlisting}

perché D-3.6 governa soltanto il mezzo usato dalla specifica interazione.

\subsection{classify -- la documentazione assegna un semantic kind}

\begin{examplebox}[DOCUMENTAZIONE DI PARTENZA -- DIDATTICO NON NORMATIVO]
\code{AuditArchive} è un servizio di conservazione dei record di audit.
\end{examplebox}

Se la distinzione \code{Service} è materialmente utile e governata:

\begin{lstlisting}
BAReferent:
  AuditArchive
  Service

BAProposition:
  classify
    classifiedReferent -> AuditArchive
    semanticKind       -> Service
\end{lstlisting}

\textbf{Non estrarre:} protocollo del servizio, storage technology o deployment. Una classificazione non sostituisce relazioni funzionali più specifiche.

\subsection{decisionRule -- la documentazione governa una decisione condizionale}

\begin{examplebox}[DOCUMENTAZIONE DI PARTENZA -- Facial Access FR-1.1]
\code{ControlledAreaAccess} può produrre un \code{AccessDecision} che consenta l'accesso soltanto quando l'\code{IdentityDeterminationOutcome} rappresenta una determinazione riuscita e la condizione di autorizzazione richiesta risulta soddisfatta per la stessa \code{GovernedIdentity}. Un outcome \code{NEGATIVE} o \code{INCONCLUSIVE} non può consentire l'accesso.
\end{examplebox}

La documentazione governa actor, input, result e una condizione composta.

\begin{lstlisting}
BAReferent:
  ControlledAreaAccess
  IdentityDeterminationOutcome
  AccessAuthorizationState
  RequiredAccessAuthorizationCondition
  AccessDecision

BAProposition:
  decisionRule
    actor  -> ControlledAreaAccess
    input  -> IdentityDeterminationOutcome
    input  -> AccessAuthorizationState
    result -> AccessDecision

    rule
      IF allOf(
           comparison(
             referent = IdentityDeterminationOutcome,
             property = outcomeKind,
             comparisonKey = equals,
             value = SUCCESS),
           satisfies(
             subject = AccessAuthorizationState,
             condition = RequiredAccessAuthorizationCondition))
      THEN
        resultAssignment
          target -> AccessDecision
          value  -> ALLOWING_ACCESS_DECISION
      ELSE
        resultAssignment
          target -> AccessDecision
          value  -> NON_ALLOWING_ACCESS_DECISION
\end{lstlisting}

Il punto essenziale è \code{satisfies}: la documentazione governa che una condizione di autorizzazione sia soddisfatta, ma non governa un property model interno come:

\begin{lstlisting}
AccessAuthorizationState.authorized = TRUE
\end{lstlisting}

Quindi quel boolean non viene inventato.

\section{Dizionario operativo dei roleKey}

I roleKey dicono \emph{che funzione svolge un BAReferent dentro una specifica proposition}. Non sono proprietà universali del referent. Lo stesso referent può avere ruoli diversi in proposition diverse.

\begin{longtable}{L{0.28\textwidth}L{0.25\textwidth}L{0.33\textwidth}}
\toprule
\textbf{roleKey} & \textbf{Significato} & \textbf{Come riconoscerlo nella documentazione} \\
\midrule\endhead
\code{behavior} & identità riusabile dello specifico transfer & il comportamento deve essere qualificato o target di altre proposition; Facial Access: delivery della capture. \\
\code{source} & origine del content trasferito & ``X consegna/invia ...''. \\
\code{destination} & destinatario del content & ``... a Y''. \\
\code{content} & significato trasferito & ``consegna Z''. \\
\code{actor} & soggetto che esegue l'azione & capability/component/responsibility che ``produce'', ``crea'', ``osserva'', decide o effettua la transition. \\
\code{result} & risultato reso disponibile/prodotto & ``produce/rende disponibile un esito''. \\
\code{input} & significato usato dall'operazione & ``la decisione usa X''; non aggiungerlo solo perché un dato esiste altrove. \\
\code{created} & nuovo item/event stabilito da create & ``crea un nuovo X''. \\
\code{observed} & significato esistente letto/ispezionato & ``legge/ispeziona X'' senza state change. \\
\code{subject} & elemento interessato da transition o satisfies & ciò il cui stato cambia, oppure ciò che deve soddisfare una condition. \\
\code{fromState} & stato precedente governato & ``passa da X ...''; optional se la source non governa lo stato iniziale. \\
\code{toState} & stato risultante governato & ``... allo stato Y''. \\
\code{correlatedItem} & elemento che deve mantenere binding di contesto & capture/outcome/item che deve appartenere alla stessa request/identity/evaluation. \\
\code{correlationContext} & identità che definisce il contesto condiviso & ``per la stessa request'', ``per la stessa identity''. \\
\code{referencer} & elemento che riguarda/richiama un altro & ``lo stato riguarda l'identità''. \\
\code{referenced} & significato richiamato & l'identità/oggetto a cui il referencer si riferisce. \\
\code{dependent} & significato che richiede un prerequisite & attività/capability che non può completarsi senza altro significato governato. \\
\code{prerequisite} & prerequisito richiesto & interazione, capability o condizione necessaria. \\
\code{consumer} & utilizzatore di un servizio & ``X consuma/usa il servizio Y''. \\
\code{service} & servizio consumato & servizio reso disponibile al consumer. \\
\code{provider} & provider governato del servizio & usarlo solo se la source identifica chi fornisce il servizio. \\
\code{abstract} & significato astratto da realizzare & capability/contract/behavior astratto. \\
\code{realization} & elemento concreto che materializza l'abstract & ``X è realizzato da Y''. \\
\code{responsibleParty} & parte a cui viene assegnata responsabilità & ``X è responsabile di ...''. \\
\code{responsibilityScope} & oggetto/ambito della responsabilità & ciò di cui la parte è responsabile. \\
\code{responsibilityKind} & tipo governato di responsabilità & approval, operation, maintenance, ecc. solo quando distinto dalla source. \\
\code{constraintTarget} & significato a cui si applica il vincolo & interaction/outcome/capability qualificato dal requisito. \\
\code{constraintValue} & valore o dominio del vincolo & valore governato, eventualmente con property key locale. \\
\code{classifiedReferent} & elemento classificato & il referent a cui viene assegnato un semantic kind. \\
\code{semanticKind} & categoria method-neutral governata & Service, Information, ecc. solo se la classificazione è materialmente utile. \\
\bottomrule
\end{longtable}

\section{Parole chiave interne a decisionRule}

Anche la struttura locale di \code{decisionRule} deve essere letta dalla documentazione, non riempita per convenienza.

\begin{longtable}{L{0.26\textwidth}L{0.62\textwidth}}
\toprule
\textbf{keyword} & \textbf{Uso operativo} \\
\midrule\endhead
\code{rule} & contiene il mapping condizione/risultato governato dalla decisione. \\
\code{decisionCondition} & condizione che la source dichiara rilevante per il risultato. \\
\code{comparison} & confronto su una property esplicitamente governata. \\
\code{property} & semantic key della proprietà confrontata; obbligatoria in comparison. \\
\code{comparisonKey} & operazione di confronto ammessa, attualmente \code{equals} o \code{notEquals}. \\
\code{value} & valore tipizzato o BAReferent con cui si confronta la property. \\
\code{satisfies} & forma locale usata quando la source governa ``X soddisfa la condizione Y'' senza governare il modello property/value interno di X. \\
\code{allOf} & tutte le condizioni figlie devono valere. \\
\code{anyOf} & almeno una condizione figlia deve valere. \\
\code{not} & negazione locale di una condition ammessa. \\
\code{resultAssignment} & assegnazione del risultato governato nel ramo THEN/ELSE. \\
\bottomrule
\end{longtable}

\begin{closedbox}[Stopping rule durante l'estrazione]
Se per completare una proposition dobbiamo inventare una property, un valore, un ruolo, un actor, uno stato o una condizione non governata, non ``completiamo'' la BA. Ci fermiamo, manteniamo \code{NOT SPECIFIED} oppure registriamo una diagnosis/counterexample se la distinzione è necessaria.
\end{closedbox}

\section{BA2 R3 - operatori}

La revisione R3 contiene quattordici operatori.

\begin{longtable}{L{0.22\textwidth}L{0.66\textwidth}}
\toprule
\textbf{Operator} & \textbf{Significato} \\
\midrule\endhead
\code{transfer} & conveyance di content da source a destination; optional behavior identity quando il trasferimento deve essere riusato/qualificato. \\
\code{produce} & actor/capability rende disponibile un risultato governato. \\
\code{create} & stabilisce un nuovo item/event project-semantic. \\
\code{observe} & legge/ispeziona significato esistente senza asserire creazione o state change. \\
\code{transition} & cambia stato/lifecycle di un subject. \\
\code{correlate} & preserva same-request/evaluation/identity/context binding. \\
\code{reference} & riferimento esplicito senza implicare dependency/correlation/ownership. \\
\code{dependOn} & prerequisite o dipendenza. \\
\code{consumeService} & consumo effettivo di servizio senza ownership transfer. \\
\code{realize} & significato concreto realizza/materializza significato astratto. \\
\code{assignResponsibility} & placement o negazione di responsibility/authority kind. \\
\code{constrain} & restriction riusabile/queryable. \\
\code{classify} & semantic kind method-neutral. \\
\code{decisionRule} & mapping operativo da condizioni/input a result. \\
\bottomrule
\end{longtable}

\section{Ruoli e cardinalità principali}

\begin{longtable}{L{0.19\textwidth}L{0.33\textwidth}L{0.34\textwidth}}
\toprule
\textbf{Operator} & \textbf{Required} & \textbf{Optional} \\
\midrule\endhead
transfer & source 1; destination 1..*; content 1..* & behavior 0..1 \\
produce & actor 1; result 1..* & input 0..* \\
create & actor 1; created 1..* & - \\
observe & actor 1; observed 1..* & result 0..* \\
transition & subject 1; toState 1 & actor 0..1; fromState 0..1 \\
correlate & correlatedItem 1..*; correlationContext 1 & - \\
reference & referencer 1; referenced 1..* & - \\
dependOn & dependent 1; prerequisite 1..* & - \\
consumeService & consumer 1..*; service 1 & provider 0..1 \\
realize & abstract 1; realization 1..* & - \\
assignResponsibility & responsibleParty 1..*; responsibilityScope 1; responsibilityKind 1 & - \\
constrain & constraintTarget 1; constraintValue 1..* & - \\
classify & classifiedReferent 1; semanticKind 1..* & - \\
decisionRule & actor 1; input 1..*; result 1 & operator-local rule \\
\bottomrule
\end{longtable}

\section{Il delta BA2 R3}

La pressione Facial Access ha giustificato soltanto due raffinamenti.

\subsection{transfer.behavior}

\begin{lstlisting}
transfer
  behavior    -> BAReferent [0..1]
  source      -> BAReferent [1]
  destination -> BAReferent [1..*]
  content     -> BAReferent [1..*]
\end{lstlisting}

Usare \code{behavior} solo quando il comportamento di trasferimento deve essere target di altre proposition.

\begin{examplebox}
\begin{lstlisting}
transfer
  behavior    -> RecognitionCaptureDelivery
  source      -> CameraSubsystem
  destination -> RecognitionProcessor
  content     -> RecognitionCapture
\end{lstlisting}

\code{RecognitionCaptureDelivery} mantiene sullo stesso segmento il consumed service, il medium corrente e le proprietà Confidentiality, Integrity e AuthorizedProvenance senza trasferirle al contenuto, agli endpoint o all'intera pipeline.
\end{examplebox}

Non si generalizza \code{behavior} agli altri operatori finché un altro counterexample non lo richiede.

\subsection{decisionRule.satisfies}

Forma minima di \code{decisionRule}:

\begin{lstlisting}
decisionRule
  actor  -> BAReferent
  input  -> BAReferent [1..*]
  result -> BAReferent

  rule
    IF    decisionCondition
    THEN  resultAssignment [1..*]
    ELSE  resultAssignment [0..*]
\end{lstlisting}

Condition forms ammesse:

\begin{lstlisting}
comparison
satisfies
allOf
anyOf
not
\end{lstlisting}

\code{comparison.property} resta obbligatoria:

\begin{lstlisting}
comparison
  referent      -> BAReferent
  property      -> controlled semantic key
  comparisonKey -> equals | notEquals
  value         -> typed local value | BAReferent
\end{lstlisting}

La forma locale \code{satisfies} è:

\begin{lstlisting}
satisfies
  subject   -> BAReferent
  condition -> BAReferent
\end{lstlisting}

Serve quando la source governa una condizione ma non ne governa il modello interno property/value.

\begin{examplebox}[FR-1.1]
\begin{lstlisting}
allOf
  comparison
    IdentityDeterminationOutcome.outcomeKind == SUCCESS

  satisfies
    subject   -> AccessAuthorizationState
    condition -> RequiredAccessAuthorizationCondition
\end{lstlisting}

Questo preserva ``authorization condition satisfied'' senza inventare:
\begin{lstlisting}
AccessAuthorizationState.authorized = TRUE
\end{lstlisting}
\end{examplebox}

\subsection{Cosa R3 non introduce}

\begin{lstlisting}
acquire operator
behavior role generalized to all operators
property-less comparison
AccessAuthorizationState.authorized
TRUE/FALSE authorization vocabulary
general-purpose predicate DSL
top-level satisfies operator
\end{lstlisting}

\section{constrain e property vocabulary}

Quando la documentazione governa un dominio di valori ma non la regola di selezione:

\begin{lstlisting}
constrain
  constraintTarget -> IdentityDeterminationOutcome
  constraintValue
    property   -> outcomeKind
    vocabulary -> [SUCCESS, NEGATIVE, INCONCLUSIVE]
\end{lstlisting}

Non usare \code{constrain} come escape hatch per prose non strutturate.

\section{correlate e binding di contesto}

\code{correlate} preserva binding di request, evaluation, operation o identity context.

Esempi:

\begin{lstlisting}
RecognitionCapture
    <-> IdentityDeterminationRequest

IdentityDeterminationOutcome
AccessAuthorizationState
    <-> same GovernedIdentity

IdentityDeterminationOutcome
    <-> AccessAttempt
\end{lstlisting}

Non usare \code{reference} come sostituto quando il requisito richiede identità di contesto.

\section{Quando non scegliere un operatore}

Il verbo documentale \code{acquire} non forza un nuovo operatore.

\begin{lstlisting}
create
observe
produce
combination
\end{lstlisting}

sono interpretazioni materialmente diverse.

La BA Facial Access conserva:

\begin{lstlisting}
RecognitionCaptureAcquisition [BAReferent]
diagnosis:
    MULTIPLE MATERIAL BA CANDIDATES
\end{lstlisting}

e non materializza una proposition falsa soltanto per ottenere un arco nel modello.

\section{BA3 - provenance, derivazione e change}

Ogni BA element meaning-bearing deve risolvere almeno a una baseline governata e alla provenance richiesta dal contratto BA3.

Origin states:

\begin{lstlisting}
GROUNDED
DERIVED
DIAGNOSTIC_UNRESOLVED
\end{lstlisting}

Tre responsabilità restano distinte:

\begin{lstlisting}
sourceLink
derivationBasisBinding
revalidationContext
\end{lstlisting}

Non usare provenance per simulare una relazione che BA2 non sa rappresentare.

\subsection{Derivazione}

Un elemento \code{DERIVED} richiede role-bound derivation basis, immutable derivation-rule revision, applicability contract, output semantic contract e lineage fino a governed source.

Nel baseline Facial Access corrente le proposition accettate sono normalizzazioni GROUNDED; non è necessario inventare derived semantics per completare la BA.

\subsection{Change e revalidation}

Quando la source cambia:

\begin{lstlisting}
source change
    -> explicit impact / revalidation context
    -> impacted BA elements
    -> PENDING_REVIEW / STALE where applicable
    -> RETAIN / REPLACE / RETIRE
    -> accepted BA for new baseline
    -> rebuild projections
\end{lstlisting}

Una baseline BA accettata non viene silenziosamente mutata e riutilizzata come se le analisi storiche avessero usato la nuova semantica.

\section{Diagnostics e NOT SPECIFIED}

La BA può accettare una diagnosis come diagnosis.

Vocabolario operativo utile:

\begin{lstlisting}
MULTIPLE MATERIAL BA CANDIDATES
BA REQUIRES UNSUPPORTED FACT
BA LOSES GOVERNED MATERIAL DISTINCTION
BA EXPOSES SOURCE CONFLICT
NOT SPECIFIED
\end{lstlisting}

\begin{closedbox}
Una diagnosis non è una project truth. \code{NOT SPECIFIED} non significa false, denied o absent; significa che il significato necessario non è governato dalla baseline corrente.
\end{closedbox}

\section{BA5 - authoring canonico}

I semantic keys operativi devono essere canonici e risolti nel corretto semantic domain.

Domini distinti:

\begin{itemize}
\item project referent names;
\item BA2 operator keys;
\item operator-scoped role keys;
\item semantic kind keys;
\item controlled value domains.
\end{itemize}

Un synonym, una traduzione o una similarità lessicale non stabiliscono equivalenza normativa automaticamente.

\subsection{Registry interpretation corrente}

BA5 non può cambiare la semantica di un operatore o di un role contract. Per operatori e ruoli la registry authority corrente è BA2 R3.

Lo snapshot a tredici operatori incluso nel corpo storico BA5 R1 è closure-time evidence e non deve essere usato come current operator list.

\section{BA4 - projection boundary}

Una projection è downstream.

\begin{lstlisting}
governed documentation
    -> accepted BA
        -> projection
            -> method-owned interpretation
\end{lstlisting}

Ogni meaning-bearing projection item deve tracciare alla BA.

Coverage modes:

\begin{lstlisting}
EXHAUSTIVE_FOR_DECLARED_SCOPE
SELECTIVE
\end{lstlisting}

Una selective projection può omettere elementi BA senza dichiararli falsi o assenti. Una method-owned interpretation può essere più ricca della shared rendering, ma resta downstream, trace-bound e non diventa BA/project authority.

\section{Facial Access - regression-passed case evidence}

La baseline analitica del caso è:

\begin{lstlisting}
FACIAL-ACCESS-BA-R24-R1
\end{lstlisting}

con project authority:

\begin{lstlisting}
FACIAL-ACCESS-GOV-R2
\end{lstlisting}

La source revision usata per la costruzione/regression del caso è:

\begin{lstlisting}
8af2257a1df94fa5a83d4853ed0a1eb4d020c429
\end{lstlisting}

Questo pin appartiene all'evidenza Facial Access e non va confuso con il continuation baseline R25 della presente guida.

\subsection{Referent principali}

\begin{lstlisting}
ControlledAreaAccess
AccessAttempt
AccessAuthorizationManagement
IdentityDetermination
PersonPresentAtAccessPoint
IdentityDeterminationRequest
RecognitionCapture
RecognitionCaptureAcquisition
RecognitionCaptureDelivery
CameraSubsystem
RecognitionProcessor
IdentityDeterminationOutcome
GovernedIdentity
AccessAuthorizationState
RequiredAccessAuthorizationCondition
AccessDecision
ConnectivityService
Confidentiality
Integrity
AuthorizedProvenance
\end{lstlisting}

\subsection{Flow minimo}

\begin{lstlisting}
PersonPresentAtAccessPoint
    -> IdentityDeterminationRequest

CameraSubsystem
    -> RecognitionCapture
       bound to IdentityDeterminationRequest

RecognitionCaptureDelivery
    CameraSubsystem -> RecognitionProcessor

RecognitionProcessor
    -> IdentityDeterminationOutcome
       SUCCESS | NEGATIVE | INCONCLUSIVE

SUCCESS
    -> GovernedIdentity X

AccessAuthorizationManagement
    -> AccessAuthorizationState concerning X

ControlledAreaAccess
    -> AccessDecision for AccessAttempt
\end{lstlisting}

\subsection{Security scope}

\begin{lstlisting}
RecognitionCaptureDelivery
    -> Confidentiality
    -> Integrity
    -> AuthorizedProvenance
\end{lstlisting}

Non propagare automaticamente quelle proprietà a outcome, authorization state, access decision o intero sistema.

\subsection{Worked propositions}

\begin{lstlisting}
P-07
correlate
  correlatedItem     -> RecognitionCapture
  correlationContext -> IdentityDeterminationRequest

P-08
transfer
  behavior    -> RecognitionCaptureDelivery
  source      -> CameraSubsystem
  destination -> RecognitionProcessor
  content     -> RecognitionCapture

P-09
consumeService
  consumer -> IdentityDeterminationRealization
  service  -> ConnectivityService

P-12
constrain
  constraintTarget -> RecognitionCaptureDelivery
  constraintValue
    property   -> interconnectionMedium
    vocabulary -> [WIRED_ETHERNET]

P-17
reference
  referencer -> AccessAuthorizationState
  referenced -> GovernedIdentity

P-23
decisionRule
  actor  -> ControlledAreaAccess
  input  -> IdentityDeterminationOutcome
  input  -> AccessAuthorizationState
  result -> AccessDecision
\end{lstlisting}

La lista completa, le condizioni di P-23 e la source/provenance matrix restano nell'artifact \code{R24_FACIAL_ACCESS_BASE_ANALYSIS_R1.md}.

\section{Controlled usefulness test}

La BA governa:

\begin{lstlisting}
RecognitionCaptureDelivery.interconnectionMedium
    = WIRED_ETHERNET
\end{lstlisting}

Un consumer può concludere che una precondizione necessaria ``this governed delivery uses RF/wireless medium'' non è soddisfatta dalla interaction governata corrente.

Non può inferire:

\begin{lstlisting}
the consumed provider infrastructure
is wired end-to-end
\end{lstlisting}

Questa distinzione mostra perché la BA deve preservare abbastanza semantica per domande downstream senza inventare topologia o provider-internal meaning.

\section{Regression procedure}

Dopo una material BA change o un contract refinement:

\begin{enumerate}
\item ricostruire/aggiornare la BA contro la source pin;
\item verificare ogni proposition materiale;
\item verificare diagnostics e forbidden inferences;
\item controllare property scope;
\item controllare cross-MR identity/request binding;
\item rieseguire usefulness/projection checks appropriati;
\item cercare un nuovo counterexample;
\item riaprire solo il contratto minimo se necessario.
\end{enumerate}

La full Facial Access post-BA regression ha confermato BA2 R3 senza richiedere reopen di BA1, BA3, BA4 o BA5 e senza trovare nuova BA2 pressure.

\section{BA6 - integrated validation gate ancora aperto}

La precedente evidenza Facial Access dimostra una parte utile dell'integrazione:

\begin{lstlisting}
complete Facial Access BA
source-to-BA coverage
post-BA regression PASS
provenance/source drill-down
controlled projection-readiness check
no new Facial Access BA2 pressure
\end{lstlisting}

Non dimostra ancora l'intero gate BA6 ereditato.

Restano richiesti:

\begin{lstlisting}
structurally different holdout
multiple BA4 projections
governed change B0 -> B1
BA3 impact / revalidation
BA rebuild
projection rebuild / re-analysis
integrated regression verdict
\end{lstlisting}

\begin{openbox}[BA6 status]
\begin{lstlisting}
BA6 integrated acceptance
    OPEN
\end{lstlisting}
La precedente closure attempt resta historical/provisional evidence e non autorizza la claim ``BA6 PASS/CLOSED''.
\end{openbox}

\section{Cosa non è la Base Analysis}

La BA non è:

\begin{itemize}
\item un DFD obbligatorio;
\item un modello di deployment completo;
\item una tassonomia STRIDE;
\item una sostituzione di SysML/ArchiMate/AADL;
\item una inference engine che completa il progetto;
\item una raccolta di security mechanisms;
\item la project authority;
\item un database/grafo obbligatorio.
\end{itemize}

\section{Checklist operativa}

Prima di accettare una BA baseline chiedere:

\begin{enumerate}
\item baseline authority e source revision sono pinned?
\item sto usando il current pre-holdout contract set corretto?
\item ogni referent riusabile ha identity stabile?
\item ogni proposition usa operator/role canonici da BA2 R3?
\item sto usando una proposition come proxy di un comportamento?
\item ho introdotto un fatto che la source non governa?
\item un \code{NOT SPECIFIED} è stato trasformato in default?
\item ogni elemento ha provenance sufficiente?
\item le proprietà locali restano scoped al target corretto?
\item cross-request / cross-identity bindings sono preservati?
\item la regression completa per lo scope dichiarato passa?
\item una projection controllata consuma la BA senza reparsing del prose per shared meaning già rappresentato?
\item esiste un counterexample corrente che richiede reopen?
\item sto confondendo case regression PASS con integrated BA6 acceptance?
\end{enumerate}

\section{Stato R25 corrente}

\begin{lstlisting}
pre-holdout BA contract
    BA0 R1
    BA1 R1
    BA2 R3
    BA3 R1
    BA4 R1
    BA5 R1

Facial Access case BA
    FACIAL-ACCESS-BA-R24-R1
    complete case BA
    post-BA regression PASS

BA6 integrated acceptance
    OPEN

holdout
    NOT SELECTED
    NOT AUTHORED

ThreatForge
    DEFERRED

STRIDE / STRIDE-AI
    NOT STARTED
\end{lstlisting}

Questa guida R3 è la guida human-readable operativa congelata per il Base Analysis pre-holdout contract di R25 Phase 2. Il checkpoint Phase 2 ne registra l'identità insieme al contract set e agli hash finali.

Dopo Phase 2 il passo autorizzato è R25 Phase 3 - thesis-ready evidence consolidation. Questo non autorizza ancora holdout selection, substantive thesis rewrite, ThreatForge implementation o STRIDE / STRIDE-AI.

\end{document}
